\chapter*{Quick Reference}
\addcontentsline{toc}{chapter}{Quick Reference}

This appendix provide a brief listing of the Fortran language bindings of the
CDI library routines:

\section*{\tt \htmlref{gridDefNvertex}{gridDefNvertex}}

\begin{verbatim}
    SUBROUTINE gridDefNvertex (INTEGER gridID, INTEGER nvertex)
\end{verbatim}

Define the number of vertex of a Gridbox.


\section*{\tt \htmlref{gridDefXbounds}{gridDefXbounds}}

\begin{verbatim}
    SUBROUTINE gridDefXbounds (INTEGER gridID, DOUBLEPRECISION xbounds(*))
\end{verbatim}

Define the bounds of a X-axis.


\section*{\tt \htmlref{gridDefYbounds}{gridDefYbounds}}

\begin{verbatim}
    SUBROUTINE gridDefYbounds (INTEGER gridID, DOUBLEPRECISION ybounds(*))
\end{verbatim}

Define the bounds of a Y-axis.


\section*{\tt \htmlref{gridInqNvertex}{gridInqNvertex}}

\begin{verbatim}
    INTEGER FUNCTION gridInqNvertex (INTEGER gridID)
\end{verbatim}

Get the number of vertex of a Gridbox.


\section*{\tt \htmlref{gridInqUUID}{gridInqUUID}}

\begin{verbatim}
    SUBROUTINE gridInqUUID (INTEGER gridID, INTEGER*1(16) uuid)
\end{verbatim}

Get the UUID of an unstructured grid.


\section*{\tt \htmlref{gridInqXbounds}{gridInqXbounds}}

\begin{verbatim}
    INTEGER FUNCTION gridInqXbounds (INTEGER gridID, DOUBLEPRECISION xbounds(*))
\end{verbatim}

Get the bounds of a X-axis.


\section*{\tt \htmlref{gridInqYbounds}{gridInqYbounds}}

\begin{verbatim}
    INTEGER FUNCTION gridInqYbounds (INTEGER gridID, DOUBLEPRECISION ybounds(*))
\end{verbatim}

Get the bounds of a Y-axis.


\section*{\tt \htmlref{subtypeCreate}{subtypeCreate}}

\begin{verbatim}
    INTEGER FUNCTION subtypeCreate (INTEGER subtype)
\end{verbatim}

Create a variable subtype.


\section*{\tt \htmlref{subtypeDefActiveIndex}{subtypeDefActiveIndex}}

\begin{verbatim}
    SUBROUTINE subtypeDefActiveIndex (INTEGER subtypeID, INTEGER index)
\end{verbatim}

Set the currently active index of a subtype (e.g. current tile index)..


\section*{\tt \htmlref{subtypeInqActiveIndex}{subtypeInqActiveIndex}}

\begin{verbatim}
    INTEGER FUNCTION subtypeInqActiveIndex (INTEGER subtypeID)
\end{verbatim}

Get the currently active index of a subtype (e.g. current tile index)..


\section*{\tt \htmlref{subtypeInqSize}{subtypeInqSize}}

\begin{verbatim}
    INTEGER FUNCTION subtypeInqSize (INTEGER subtypeID)
\end{verbatim}

Get the size of a subtype (e.g. no. of tiles)..


\section*{\tt \htmlref{subtypeInqTile}{subtypeInqTile}}

\begin{verbatim}
    INTEGER FUNCTION subtypeInqTile (INTEGER subtypeID, INTEGER tileindex,
                                     INTEGER attribute)
\end{verbatim}

Specialized version of subtypeInqSubEntry looking for tile/attribute pair..


\section*{\tt \htmlref{tableRead}{tableRead}}

\begin{verbatim}
    INTEGER FUNCTION tableRead (CHARACTER*(*) tablefile)
\end{verbatim}

read table of parameters from file in tabular format.


\section*{\tt \htmlref{tableWrite}{tableWrite}}

\begin{verbatim}
    SUBROUTINE tableWrite (CHARACTER*(*) filename, INTEGER tableID)
\end{verbatim}

write table of parameters to file in tabular format.


\section*{\tt \htmlref{tableWriteC}{tableWriteC}}

\begin{verbatim}
    SUBROUTINE tableWriteC (CHARACTER*(*) filename, INTEGER tableID)
\end{verbatim}

write table of parameters to file in C language format.


\section*{\tt \htmlref{taxisCreate}{taxisCreate}}

\begin{verbatim}
    INTEGER FUNCTION taxisCreate (INTEGER timetype)
\end{verbatim}

Create a Time axis.


\section*{\tt \htmlref{taxisDefCalendar}{taxisDefCalendar}}

\begin{verbatim}
    SUBROUTINE taxisDefCalendar (INTEGER taxisID, INTEGER calendar)
\end{verbatim}

Define the calendar.


\section*{\tt \htmlref{taxisDefFdate}{taxisDefFdate}}

\begin{verbatim}
    SUBROUTINE taxisDefFdate (INTEGER taxisID, INTEGER date)
\end{verbatim}

Define the forecast reference date.


\section*{\tt \htmlref{taxisDefFtime}{taxisDefFtime}}

\begin{verbatim}
    SUBROUTINE taxisDefFtime (INTEGER taxisID, INTEGER time)
\end{verbatim}

Define the forecast reference time.


\section*{\tt \htmlref{taxisDefRdate}{taxisDefRdate}}

\begin{verbatim}
    SUBROUTINE taxisDefRdate (INTEGER taxisID, INTEGER date)
\end{verbatim}

Define the reference date.


\section*{\tt \htmlref{taxisDefRtime}{taxisDefRtime}}

\begin{verbatim}
    SUBROUTINE taxisDefRtime (INTEGER taxisID, INTEGER time)
\end{verbatim}

Define the reference time.


\section*{\tt \htmlref{taxisDefVdate}{taxisDefVdate}}

\begin{verbatim}
    SUBROUTINE taxisDefVdate (INTEGER taxisID, INTEGER date)
\end{verbatim}

Define the verification date.


\section*{\tt \htmlref{taxisDefVtime}{taxisDefVtime}}

\begin{verbatim}
    SUBROUTINE taxisDefVtime (INTEGER taxisID, INTEGER time)
\end{verbatim}

Define the verification time.


\section*{\tt \htmlref{taxisDestroy}{taxisDestroy}}

\begin{verbatim}
    SUBROUTINE taxisDestroy (INTEGER taxisID)
\end{verbatim}

Destroy a Time axis.


\section*{\tt \htmlref{taxisInqCalendar}{taxisInqCalendar}}

\begin{verbatim}
    INTEGER FUNCTION taxisInqCalendar (INTEGER taxisID)
\end{verbatim}

Get the calendar.


\section*{\tt \htmlref{taxisInqFdate}{taxisInqFdate}}

\begin{verbatim}
    INTEGER FUNCTION taxisInqFdate (INTEGER taxisID)
\end{verbatim}

Get the forecast reference date.


\section*{\tt \htmlref{taxisInqFtime}{taxisInqFtime}}

\begin{verbatim}
    INTEGER FUNCTION taxisInqFtime (INTEGER taxisID)
\end{verbatim}

Get the forecast reference time.


\section*{\tt \htmlref{taxisInqRdate}{taxisInqRdate}}

\begin{verbatim}
    INTEGER FUNCTION taxisInqRdate (INTEGER taxisID)
\end{verbatim}

Get the reference date.


\section*{\tt \htmlref{taxisInqRtime}{taxisInqRtime}}

\begin{verbatim}
    INTEGER FUNCTION taxisInqRtime (INTEGER taxisID)
\end{verbatim}

Get the reference time.


\section*{\tt \htmlref{taxisInqVdate}{taxisInqVdate}}

\begin{verbatim}
    INTEGER FUNCTION taxisInqVdate (INTEGER taxisID)
\end{verbatim}

Get the verification date.


\section*{\tt \htmlref{taxisInqVtime}{taxisInqVtime}}

\begin{verbatim}
    INTEGER FUNCTION taxisInqVtime (INTEGER taxisID)
\end{verbatim}

Get the verification time.


\section*{\tt \htmlref{vlistInqVarSubtype}{vlistInqVarSubtype}}

\begin{verbatim}
    INTEGER FUNCTION vlistInqVarSubtype (INTEGER vlistID, INTEGER varID)
\end{verbatim}

Return subtype ID for a given variable..


\section*{\tt \htmlref{zaxisCreate}{zaxisCreate}}

\begin{verbatim}
    INTEGER FUNCTION zaxisCreate (INTEGER zaxistype, INTEGER size)
\end{verbatim}

Create a vertical Z-axis.


\section*{\tt \htmlref{zaxisDefLevel}{zaxisDefLevel}}

\begin{verbatim}
    SUBROUTINE zaxisDefLevel (INTEGER zaxisID, INTEGER levelID,
                              DOUBLEPRECISION levels)
\end{verbatim}

Define one level of a Z-axis.


\section*{\tt \htmlref{zaxisDefLevels}{zaxisDefLevels}}

\begin{verbatim}
    SUBROUTINE zaxisDefLevels (INTEGER zaxisID, DOUBLEPRECISION levels(*))
\end{verbatim}

Define the levels of a Z-axis.


\section*{\tt \htmlref{zaxisDefLongname}{zaxisDefLongname}}

\begin{verbatim}
    SUBROUTINE zaxisDefLongname (INTEGER zaxisID, CHARACTER*(*) longname_optional)
\end{verbatim}

Define the longname of a Z-axis.


\section*{\tt \htmlref{zaxisDefName}{zaxisDefName}}

\begin{verbatim}
    SUBROUTINE zaxisDefName (INTEGER zaxisID, CHARACTER*(*) name_optional)
\end{verbatim}

Define the name of a Z-axis.


\section*{\tt \htmlref{zaxisDefNlevRef}{zaxisDefNlevRef}}

\begin{verbatim}
    SUBROUTINE zaxisDefNlevRef (INTEGER gridID, INTEGER nhlev)
\end{verbatim}

Define the number of half levels of a generalized Z-axis.


\section*{\tt \htmlref{zaxisDefNumber}{zaxisDefNumber}}

\begin{verbatim}
    SUBROUTINE zaxisDefNumber (INTEGER gridID, INTEGER number)
\end{verbatim}

Define the reference number for a generalized Z-axis.


\section*{\tt \htmlref{zaxisDefPsName}{zaxisDefPsName}}

\begin{verbatim}
    SUBROUTINE zaxisDefPsName (INTEGER zaxisID, CHARACTER*(*) psname_optional)
\end{verbatim}

Define the name of the surface pressure variable of a hybrid sigma pressure Z-axis.


\section*{\tt \htmlref{zaxisDefUUID}{zaxisDefUUID}}

\begin{verbatim}
    SUBROUTINE zaxisDefUUID (INTEGER zaxisID, INTEGER*1(16) uuid)
\end{verbatim}

Define the UUID of a generalized Z-axis.


\section*{\tt \htmlref{zaxisDefUnits}{zaxisDefUnits}}

\begin{verbatim}
    SUBROUTINE zaxisDefUnits (INTEGER zaxisID, CHARACTER*(*) units_optional)
\end{verbatim}

Define the units of a Z-axis.


\section*{\tt \htmlref{zaxisDestroy}{zaxisDestroy}}

\begin{verbatim}
    SUBROUTINE zaxisDestroy (INTEGER zaxisID)
\end{verbatim}

Destroy a vertical Z-axis.


\section*{\tt \htmlref{zaxisDuplicate}{zaxisDuplicate}}

\begin{verbatim}
    INTEGER FUNCTION zaxisDuplicate (INTEGER zaxisID)
\end{verbatim}

Duplicate a Z-axis.


\section*{\tt \htmlref{zaxisInqLevel}{zaxisInqLevel}}

\begin{verbatim}
    DOUBLEPRECISION FUNCTION zaxisInqLevel (INTEGER zaxisID, INTEGER levelID)
\end{verbatim}

Get one level of a Z-axis.


\section*{\tt \htmlref{zaxisInqLevels}{zaxisInqLevels}}

\begin{verbatim}
    SUBROUTINE zaxisInqLevels (INTEGER zaxisID, DOUBLEPRECISION levels(*))
\end{verbatim}

Get all levels of a Z-axis.


\section*{\tt \htmlref{zaxisInqLongname}{zaxisInqLongname}}

\begin{verbatim}
    SUBROUTINE zaxisInqLongname (INTEGER zaxisID, CHARACTER*(*) longname)
\end{verbatim}

Get the longname of a Z-axis.


\section*{\tt \htmlref{zaxisInqName}{zaxisInqName}}

\begin{verbatim}
    SUBROUTINE zaxisInqName (INTEGER zaxisID, CHARACTER*(*) name)
\end{verbatim}

Get the name of a Z-axis.


\section*{\tt \htmlref{zaxisInqNlevRef}{zaxisInqNlevRef}}

\begin{verbatim}
    INTEGER FUNCTION zaxisInqNlevRef (INTEGER gridID)
\end{verbatim}

Get the number of half levels of a generalized Z-axis.


\section*{\tt \htmlref{zaxisInqNumber}{zaxisInqNumber}}

\begin{verbatim}
    INTEGER FUNCTION zaxisInqNumber (INTEGER gridID)
\end{verbatim}

Get the reference number to a generalized Z-axis.


\section*{\tt \htmlref{zaxisInqPsName}{zaxisInqPsName}}

\begin{verbatim}
    SUBROUTINE zaxisInqPsName (INTEGER zaxisID, CHARACTER*(*) psname)
\end{verbatim}

Get the name of the surface pressure variable of a hybrid sigma pressure Z-axis.


\section*{\tt \htmlref{zaxisInqSize}{zaxisInqSize}}

\begin{verbatim}
    INTEGER FUNCTION zaxisInqSize (INTEGER zaxisID)
\end{verbatim}

Get the size of a Z-axis.


\section*{\tt \htmlref{zaxisInqStdname}{zaxisInqStdname}}

\begin{verbatim}
    SUBROUTINE zaxisInqStdname (INTEGER zaxisID, CHARACTER*(*) stdname)
\end{verbatim}

Get the standard name of a Z-axis.


\section*{\tt \htmlref{zaxisInqType}{zaxisInqType}}

\begin{verbatim}
    INTEGER FUNCTION zaxisInqType (INTEGER zaxisID)
\end{verbatim}

Get the type of a Z-axis.


\section*{\tt \htmlref{zaxisInqUUID}{zaxisInqUUID}}

\begin{verbatim}
    SUBROUTINE zaxisInqUUID (INTEGER zaxisID, INTEGER*1(16) uuid)
\end{verbatim}

Get the UUID of a generalized Z-axis.


\section*{\tt \htmlref{zaxisInqUnits}{zaxisInqUnits}}

\begin{verbatim}
    SUBROUTINE zaxisInqUnits (INTEGER zaxisID, CHARACTER*(*) units)
\end{verbatim}

Get the units of a Z-axis.


